\chapter{Geometría normada y dimensión finita}\label{cap:nuevo08}\label{chap:19}\label{chap:22}

Un espacio normado reúne dos estructuras: podemos sumar y multiplicar vectores, y también medir distancias. Comenzaremos mostrando que ambas estructuras son compatibles; después estudiaremos series, subespacios cerrados y direcciones que permanecen alejadas de un subespacio propio. Estas herramientas permitirán reconocer qué fenómenos dependen realmente de la dimensión finita.

En $\R^n$ todas las normas describen la misma convergencia, toda transformación lineal es continua y los conjuntos cerrados acotados son compactos. No son propiedades de una norma aislada: el número finito de coordenadas es decisivo. Explicaremos tanto estas afirmaciones como las obstrucciones que aparecen en dimensión infinita.

Trabajaremos sobre $\F=\R$ o $\C$. Se requieren completitud, separabilidad, compacidad y el teorema de Baire. Una dimensión sin calificativo será algebraica: una base de Hamel permite únicamente combinaciones lineales finitas. No se presupone la teoría de operadores del capítulo siguiente; las afirmaciones sobre aplicaciones lineales en dimensión finita se demostrarán directamente mediante coordenadas.

\section{Operaciones, bolas y coordenadas}\label{lectura:19}

\begin{proposicion}\label{ch19:operaciones}
En un espacio normado $V$, la suma $V\times V\to V$ y el producto $\F\times V\to V$ son continuos. La norma es una función de Lipschitz de constante $1$.
\end{proposicion}
\begin{proof}
La desigualdad triangular, aplicada en ambos sentidos, da
\[
 \big|\norm{x}-\norm{y}\big|\leq\norm{x-y}.
\]
Para la suma, fijados $x_0,y_0$, tenemos
\[
 \norm{(x+y)-(x_0+y_0)}\leq\norm{x-x_0}+\norm{y-y_0}.
\]
Si ambos sumandos del lado derecho son menores que $\eps/2$, el resultado es menor que $\eps$. Para el producto, separamos los dos cambios mediante la identidad
\[
 \alpha x-\alpha_0x_0=\alpha(x-x_0)+(\alpha-\alpha_0)x_0.
\]
Si $|\alpha-\alpha_0|<\delta\leq1$, entonces $|\alpha|<|\alpha_0|+1$. Por tanto, si además $\norm{x-x_0}<\delta$, resulta
\[
 \norm{\alpha x-\alpha_0x_0}
 <\bigl(|\alpha_0|+1+\norm{x_0}\bigr)\delta.
\]
Basta elegir $\delta=\min\{1,\eps/(|\alpha_0|+1+\norm{x_0})\}$. La elección depende del punto de continuidad, como permite la definición.
\end{proof}

\begin{corolario}\label{ch19:limites}
Si $x_n\to x$, $y_n\to y$ y $\alpha_n\to\alpha$, entonces $x_n+y_n\to x+y$ y $\alpha_nx_n\to\alpha x$.
\end{corolario}
\begin{proof}
Aplicamos las dos estimaciones anteriores a los pares de términos de las sucesiones. En la segunda, $\alpha_n$ es acotada porque converge, de modo que $|\alpha_n|\norm{x_n-x}+|\alpha_n-\alpha|\norm{x}\to0$.
\end{proof}

\begin{observacion}
En el corolario de la lectura original aparece $\alpha_n+v_n$ donde corresponde $\alpha_nv_n$. En un espacio vectorial general ni siquiera está definida la suma de un escalar con un vector.
\end{observacion}

\begin{lema}[Coordenadas en dimensión finita]\label{ch19:coordenadas}
Si $v_1,\ldots,v_m$ son linealmente independientes en $V$, existen $c,C>0$ tales que
\[
 c\max_i|a_i|\leq\norm{\sum_{i=1}^m a_iv_i}\leq C\max_i|a_i|
 \qquad(a_1,\ldots,a_m\in\F).
\]
En particular, todo subespacio de dimensión finita es completo y cerrado en $V$.
\end{lema}
\begin{proof}
La cota superior vale con $C=\sum_i\norm{v_i}$. Definamos $L(a)=\sum_i a_iv_i$. Esa misma cota da $\norm{L(a)-L(b)}\leq C\norm{a-b}_\infty$, así que $a\mapsto\norm{L(a)}$ es continua. La esfera $S=\{a\in\F^m:\norm a_\infty=1\}$ es compacta por Heine--Borel en el espacio euclidiano $\R^m$ o $\R^{2m}$. La función alcanza allí un mínimo $c$. La independencia implica $L(a)\ne0$ en $S$, y por ello $c>0$. Para $a\ne0$, aplicamos esta cota a $a/\norm a_\infty$ y multiplicamos por $\norm a_\infty$; para $a=0$ la desigualdad es inmediata.

Sea ahora $x_n=L(a^{(n)})$ una sucesión de Cauchy en el subespacio generado. La cota inferior implica que cada sucesión de coordenadas es de Cauchy en $\F$. Si $a^{(n)}\to a$, la cota superior da $x_n\to L(a)$ en el subespacio. Este es completo. Si una sucesión del subespacio converge en $V$, es de Cauchy y el límite que acabamos de obtener pertenece al subespacio; la unicidad del límite demuestra que este es cerrado. El subespacio $\{0\}$ satisface ambas conclusiones directamente.
\end{proof}

Este lema se ha demostrado antes de aplicar Baire: sus consecuencias sobre cerradura y completitud se necesitan ya en las propiedades generales que siguen. La sección~\ref{sec:nuevo08:finita} desarrollará después el papel de la dimensión finita de manera sistemática.

\begin{proposicion}\label{ch19:geometria}
Sea $V$ un espacio normado y $r>0$.
\begin{enumerate}[label=(\roman*)]
\item $\interior{\overline B(x,r)}=B(x,r)$ y $\cl{B(x,r)}=\overline B(x,r)$.
\item $V$ es conexo por caminos; dos puntos distintos pueden unirse por un arco.
\item Si $V\ne\{0\}$, entonces $V$ no es compacto.
\item $V$ es separable si y solo si contiene un subespacio denso de dimensión a lo sumo numerable.
\item Si $V$ es de Banach, su dimensión es finita o infinita no numerable.
\end{enumerate}
\end{proposicion}
\begin{proof}
(i) La continuidad de la norma muestra que la bola abierta es abierta y la cerrada es cerrada. Por ello $B(x,r)\subseteq\interior{\overline B(x,r)}$ y $\cl{B(x,r)}\subseteq\overline B(x,r)$. Si $\norm{y-x}=r$, los puntos $x+(1+s)(y-x)$, con $s>0$ tendiendo a cero, están fuera de la bola cerrada y convergen a $y$. Así, $y$ no es interior. En cambio, $x+(1-s)(y-x)$ pertenece a la bola abierta para $0<s<1$ y converge a $y$, lo que prueba la otra igualdad. Si $V=\{0\}$ no hay puntos de frontera y las igualdades siguen siendo ciertas.

(ii) Dados $u,v$, la función $\gamma(t)=(1-t)u+tv$, $0\leq t\leq1$, es continua por la proposición~\ref{ch19:operaciones}. Une $u$ con $v$. Si $u\ne v$, es inyectiva, pues $\gamma(s)=\gamma(t)$ implica $(s-t)(v-u)=0$ y, por tanto, $s=t$. Como su dominio es compacto y el codominio es Hausdorff, es un homeomorfismo sobre su imagen. La conexión por caminos implica conexión: una separación del espacio induciría una separación del intervalo a través de un camino que una puntos de las dos partes.

(iii) Elija $v\ne0$. La sucesión $nv$ tiene normas $n\norm v$ no acotadas. Por tanto $V$ no es acotado, mientras que todo compacto métrico sí lo es. La excepción $V=\{0\}$ es necesaria: ese espacio es compacto.

(iv) Si $D$ es numerable y denso, $M=\spanop D$ contiene a $D$, es denso y tiene una base a lo sumo numerable. Recíprocamente, supongamos $M=\spanop\{v_1,v_2,\ldots\}$ y $\cl M=V$. Sea $\F_0=\Q$ en el caso real y $\F_0=\Q+i\Q$ en el complejo. El conjunto $D$ de todas las combinaciones finitas de los $v_j$ con coeficientes en $\F_0$ es numerable: es una unión numerable de conjuntos numerables. Si $x=\sum_{j=1}^m a_jv_j$, elegimos $q_j\in\F_0$ tan cercanos a $a_j$ que $\sum_{j=1}^m|q_j-a_j|\norm{v_j}<\eps$. Entonces $\norm{x-\sum q_jv_j}<\eps$. Así $M\subseteq\cl D$, y al tomar cierres obtenemos $V=\cl D$.

(v) Si existiera una base infinita numerable $v_1,v_2,\ldots$, tendríamos $V=\bigcup_{n\geq1}M_n$, con $M_n=\spanop\{v_1,\ldots,v_n\}$. Cada $M_n$ es cerrado por el lema~\ref{ch19:coordenadas}. Tiene interior vacío: si $B(x,r)\subseteq M_n$, entonces $x\in M_n$, y restar $x$ da $B(0,r)\subseteq M_n$; al multiplicar por escalares se obtiene $M_n=V$, contradicción. Baire prohíbe que el espacio completo $V$ sea esa unión de cerrados con interior vacío.
\end{proof}

\index{separabilidad!espacios normados}
\index{base de Hamel}
\begin{observacion}
Separabilidad y dimensión algebraica numerable no son equivalentes. Por ejemplo, un espacio de Banach separable e infinito dimensional puede tener un subespacio denso generado por una sucesión, pero esa sucesión no constituye una base de Hamel de todo el espacio. La diferencia consiste en permitir límites además de sumas finitas.
\end{observacion}

\section{Series y completitud}

\begin{definicion}
Dada una sucesión $(v_n)$ de $V$, la serie $\sum_{n=1}^{\infty}v_n$ converge a $s$ cuando sus sumas parciales $s_N=\sum_{n=1}^N v_n$ convergen a $s$. Se llama \emph{absolutamente convergente} cuando $\sum_{n=1}^{\infty}\norm{v_n}<\infty$.
\end{definicion}
\index{serie!absolutamente convergente}

\begin{teorema}\label{ch19:series}
Un espacio normado $V$ es completo si y solo si toda serie absolutamente convergente en $V$ converge en $V$.
\end{teorema}
\begin{proof}
Si $V$ es completo y $\sum\norm{v_n}<\infty$, para $m>n$ tenemos
\[
 \norm{s_m-s_n}\leq\sum_{j=n+1}^m\norm{v_j}.
\]
Las colas de una serie real convergente de términos no negativos tienden a cero. Por ello $(s_n)$ es de Cauchy, y la completitud proporciona su límite en $V$.

Para el recíproco, sea $(x_n)$ de Cauchy. Elegimos inductivamente enteros $n_1<n_2<\cdots$ tales que
\[
 p,q\geq n_j\quad\Longrightarrow\quad\norm{x_p-x_q}<2^{-j}.
\]
La condición de Cauchy da un umbral para cada $j$; aumentarlo permite exigir además $n_{j+1}>n_j$. Definimos $v_1=x_{n_1}$ y $v_{j+1}=x_{n_{j+1}}-x_{n_j}$. Entonces
\[
 \sum_{j=1}^{\infty}\norm{v_j}
 \leq\norm{x_{n_1}}+\sum_{j=1}^{\infty}2^{-j}<\infty.
\]
Por hipótesis la serie converge a algún $x\in V$. Su suma parcial número $j$ es $x_{n_j}$, porque todos los términos intermedios se cancelan. Tenemos, pues, una subsucesión convergente.

Falta justificar que converge toda la sucesión. Dado $\eps>0$, elija $N$ tal que $\norm{x_p-x_q}<\eps/2$ para $p,q\geq N$. Después elija $j$ con $n_j\geq N$ y $\norm{x_{n_j}-x}<\eps/2$. Para todo $n\geq N$,
\[
 \norm{x_n-x}\leq\norm{x_n-x_{n_j}}+\norm{x_{n_j}-x}<\eps.
\]
Así $V$ es completo.
\end{proof}

La subsucesión anterior no se elige para simplificar una notación: convierte diferencias pequeñas en una serie numérica sumable. Es una técnica que reaparecerá en construcciones por aproximaciones sucesivas.

\section{Subespacios cerrados y lema de Riesz}

\begin{proposicion}\label{ch19:cierre-subespacio}
El cierre de un subespacio vectorial $M\subseteq V$ es un subespacio vectorial.
\end{proposicion}
\begin{proof}
Contiene al cero. Si $x,y\in\cl M$, existen $x_n,y_n\in M$ con $x_n\to x$ e $y_n\to y$. Para $\alpha,\beta\in\F$, cada $\alpha x_n+\beta y_n$ pertenece a $M$ y converge a $\alpha x+\beta y$ por el corolario~\ref{ch19:limites}. Este límite pertenece a $\cl M$.
\end{proof}

\begin{lema}[Riesz]\label{ch19:riesz}
Si $M$ es un subespacio propio cerrado de $V$, para todo $0<\theta<1$ existe $v\in V$ con $\norm v=1$ y $\dist(v,M)>\theta$. En particular, para todo $\eps>0$ puede conseguirse $\dist(v,M)\geq1-\eps$.
\end{lema}
\index{Riesz, lema de}
\begin{proof}
Tome $u\notin M$ y $d=\dist(u,M)$. El número $d$ es positivo porque $M$ es cerrado: si fuera cero, habría puntos de $M$ que convergen a $u$. Como $d/\theta>d$, la definición de ínfimo permite elegir $w\in M$ tal que $\norm{u-w}<d/\theta$. Sea $v=(u-w)/\norm{u-w}$. Para cualquier $z\in M$, el vector $w+\norm{u-w}z$ también pertenece a $M$, y por ello
\[
 \norm{v-z}
 =\frac{\norm{u-(w+\norm{u-w}z)}}{\norm{u-w}}
 \geq\frac d{\norm{u-w}}>\theta.
\]
La cota inferior es independiente de $z$, de modo que también vale para el ínfimo. Si $0<\eps<1$, tome $\theta=1-\eps$; si $\eps\geq1$, cualquier vector unitario cumple la desigualdad solicitada. Tal vector existe porque $M$ es propio.
\end{proof}

Las operaciones continuas y la caracterización de completitud mediante series describen cualquier espacio normado. El lema de Riesz añade una restricción geométrica: un subespacio cerrado propio no puede aproximar uniformemente todas las direcciones unitarias. Combinaremos ahora esa restricción con las cotas de coordenadas para caracterizar la dimensión finita.

\section{Identificar vectores: espacios cocientes}\label{ampl:cocientes}
Antes de volver a la dimensión finita estudiaremos otra operación con subespacios: declarar equivalentes dos vectores cuya diferencia pertenece a un subespacio dado. Esta construcción permite ignorar direcciones que no interesa distinguir. También resuelve una dificultad ya encontrada: una integral puede medir cero para una función que no es idénticamente nula. En ese caso hay que identificar primero los objetos que la seminorma no distingue.

La organización de esta ampliación se apoya en los apuntes y la presentación de Maximenko \cite{max-cociente,max-seminormado}. Desarrollamos las verificaciones algebraicas y los pasos con ínfimos; los ejercicios indican cuáles adaptan propuestas de esa fuente. No se requiere todavía la teoría general de operadores.
\index{espacio!cociente}\index{norma!cociente}\index{seminorma!nucleo@núcleo}

\begin{definicion}
Si $M$ es un subespacio de un espacio vectorial $V$, escribimos $x\sim y$ cuando $x-y\in M$. La clase de $x$ es $[x]=x+M$, y el conjunto de clases se denota $V/M$. Sus operaciones son
\[
 [x]+[y]=[x+y],\qquad \lambda[x]=[\lambda x].
\]
\end{definicion}
La relación es reflexiva porque $0\in M$, simétrica porque $-M=M$ y transitiva porque $M+M=M$. Las operaciones no dependen de los representantes: si $x-x'$, $y-y'\in M$, entonces $(x+y)-(x'+y')\in M$ y $\lambda x-\lambda x'\in M$. Las identidades vectoriales se heredan de $V$. El vector cero del cociente es la clase $[0]=M$, que generalmente contiene muchos vectores de $V$. La aplicación $q:V\to V/M$, $q(x)=[x]$, es lineal, sobreyectiva y tiene núcleo $M$.

\begin{proposicion}[Seminorma del cociente]\label{ampl:cociente-seminorma}
Sea $p$ una seminorma en $V$ y $M\subset V$ un subespacio. La expresión
\[
 \bar p([x])=\inf_{m\in M}p(x-m)
\]
define una seminorma en $V/M$. Es una norma si y solo si $M$ es cerrado para la pseudométrica $d_p(x,y)=p(x-y)$. En particular, si $V$ es normado y $M$ es cerrado,
\[
 \norm{[x]}_{V/M}=\dist(x,M)
\]
es una norma. La proyección $q$ satisface $\bar p(qx-qy)\le p(x-y)$.
\end{proposicion}
\begin{proof}
Los representantes de $[x]$ son exactamente $x-m$, $m\in M$; cambiar $x$ por otro representante no cambia ese conjunto. El ínfimo es finito y no negativo porque $m=0$ es una elección permitida. Para $\eps>0$ elegimos $m,n\in M$ tales que
\[
 p(x-m)<\bar p([x])+\eps/2,\qquad
 p(y-n)<\bar p([y])+\eps/2.
\]
Entonces $m+n\in M$ y
\[
 \bar p([x+y])\le p(x+y-m-n)
 <\bar p([x])+\bar p([y])+\eps.
\]
Dejar $\eps\downarrow0$ demuestra subaditividad. Para $\lambda\ne0$, al escribir $m=\lambda n$ obtenemos
\[
 \bar p([\lambda x])=\inf_{n\in M}|\lambda|p(x-n)
 =|\lambda|\bar p([x]).
\]
Para $\lambda=0$ ambos miembros son cero. Además $\bar p([x])=0$ equivale a que para todo $\eps>0$ exista $m\in M$ con $p(x-m)<\eps$, esto es, a que $x\in\cl M$ en la topología de $d_p$. Por tanto la única clase de seminorma cero es $M$ exactamente cuando $\cl M=M$. Finalmente, elegir $m=0$ en el ínfimo para $[x-y]$ prueba la última desigualdad.
\end{proof}

\begin{observacion}[El ínfimo no exige un representante óptimo]
La definición permite encontrar representantes de norma menor que $\norm{[x]}_{V/M}+\eps$ para cada $\eps>0$. No afirma que alguno tenga norma exactamente igual al ínfimo. Esta distinción será esencial en la prueba de completitud y en el ejercicio~\ref{ej:ampl-cocientes-5}.
\end{observacion}

\begin{proposicion}[Eliminar el núcleo de una seminorma]\label{ampl:nucleo-seminorma}
Sea $N_p=\{x:p(x)=0\}$. Es un subespacio cerrado para $d_p$, y
\[
 \norm{[x]}_{V/N_p}=p(x)
\]
está bien definida y es una norma.
\end{proposicion}
\begin{proof}
Si $u,v\in N_p$, entonces $p(\alpha u+\beta v)\le|\alpha|p(u)+|\beta|p(v)=0$; así $N_p$ es subespacio. La desigualdad $|p(x)-p(y)|\le p(x-y)$ prueba continuidad de $p$ respecto de $d_p$, por lo que $N_p=p^{-1}(\{0\})$ es cerrado. Si $n\in N_p$, las desigualdades triangulares en ambos sentidos dan $p(x+n)=p(x)$. Todos los representantes de una clase tienen, por tanto, el mismo valor de $p$, y su ínfimo es ese valor. La proposición anterior demuestra que la seminorma resultante es una norma.
\end{proof}

\begin{ejemplo}[Una dirección invisible]
En $\R^2$, $p(s,t)=|s|$ no distingue puntos que difieren verticalmente. Su núcleo es $\{0\}\times\R$. La aplicación $[(s,t)]\mapsto s$ identifica isométricamente el cociente con $\R$. La coordenada vertical se pierde de manera deliberada; no se está afirmando que dos pares originalmente distintos sean iguales en $\R^2$.
\end{ejemplo}

\begin{teorema}[Completitud del cociente]\label{ampl:cociente-banach}
Si $V$ es Banach y $M$ es un subespacio cerrado, $V/M$ es Banach con la norma cociente.
\end{teorema}
\begin{estrategia}
Una sucesión de clases no viene con representantes de Cauchy elegidos de antemano. Es más cómodo usar el criterio de series del teorema~\ref{ch19:series}: levantaremos una serie absolutamente convergente a otra en $V$, permitiendo un error sumable al elegir cada representante.
\end{estrategia}
\begin{proof}
Sea $\sum_{n\ge1}A_n$ una serie en $V/M$ con $\sum_n\norm{A_n}_{V/M}<\infty$. Por la propiedad del ínfimo elegimos $x_n\in A_n$ tales que
\[
 \norm{x_n}<\norm{A_n}_{V/M}+2^{-n}.
\]
Así $\sum_n\norm{x_n}<\infty$, y la completitud de $V$ da $x=\sum_nx_n\in V$. Como $q$ es lineal y $1$-Lipschitz,
\[
 \left\|\sum_{n=1}^N A_n-q(x)\right\|_{V/M}
 \le\left\|\sum_{n=1}^N x_n-x\right\|\longrightarrow0.
\]
Toda serie absolutamente convergente del cociente converge. El mismo criterio de series prueba su completitud.
\end{proof}

\begin{proposicion}[Qué ocurre con las bolas]\label{ampl:cociente-bolas}
Para $V$ normado, $M$ cerrado y $r>0$,
\[
 q(B(x,r))=B(qx,r).
\]
En particular, $q$ envía abiertos en abiertos.
\end{proposicion}
\begin{proof}
Si $y\in B(x,r)$, entonces $\norm{qy-qx}\le\norm{y-x}<r$. Recíprocamente, si $\norm{[z]-[x]}<r$, existe $m\in M$ tal que $\norm{z-x-m}<r$. El vector $z-m$ pertenece a $B(x,r)$ y representa la clase $[z]$. Para un abierto $U$ y una clase $qx\in q(U)$, con $x\in U$, existe $r>0$ tal que $B(x,r)\subset U$. La igualdad de bolas da una bola alrededor de $qx$ contenida en $q(U)$; así $q(U)$ es abierto.
\end{proof}

\section{Ejercicios sobre cocientes y seminormas}

\begin{ejercicio}[Ampliación: una recta como cociente]\label{ej:ampl-cocientes-1}
En $\R^2$ con la norma euclidiana sea $M=\{(t,t):t\in\R\}$. Calcule la norma de $[(a,b)]$, encuentre el representante de norma mínima y construya una isometría lineal de $\R^2/M$ sobre $\R$. Este ejercicio y su desarrollo son propios de esta edición.
\end{ejercicio}


\begin{ejercicio}[Ampliación: el límite como coordenada del cociente]\label{ej:ampl-cocientes-2}
Sean $c$ y $c_0$ los espacios de sucesiones escalares convergentes y convergentes a cero, con la norma uniforme. Demuestre que $c_0$ es cerrado en $c$ y que $[x]\mapsto\lim_nx_n$ es un isomorfismo isométrico de $c/c_0$ sobre $\F$. Adaptación del ejercicio de la presentación \cite{max-seminormado}; la solución se desarrolla aquí.
\end{ejercicio}


\begin{ejercicio}[Ampliación: un cociente que pierde toda distancia]\label{ej:ampl-cocientes-3}
En $\ell^2$ sea $c_{00}$ el subespacio de sucesiones de soporte finito. Pruebe que $\ell^2/c_{00}$ no es el espacio vectorial cero, aunque su seminorma cociente vale cero en todas las clases. Explique por qué esto no contradice el teorema~\ref{ampl:cociente-banach}. Adaptación del contraejemplo propuesto en \cite{max-seminormado}.
\end{ejercicio}


\begin{ejercicio}[Ampliación: observar sólo los extremos]\label{ej:ampl-cocientes-4}
En $C([0,1],\F)$ defina $p(f)=|f(0)|+|f(1)|$. Describa su núcleo y pruebe que el cociente por ese núcleo es isométrico a $\F^2$ con la norma $\norm{(a,b)}_1=|a|+|b|$. Determine si el cociente es completo aunque $p$ no sea una norma. Ejercicio propio de esta edición.
\end{ejercicio}


\begin{ejercicio}[Ampliación: norma cociente sin mínimo]\label{ej:ampl-cocientes-5}
Sea $\varphi:V\to\F$ lineal y no nula, y suponga
$C=\sup_{\norm v\le1}|\varphi(v)|<\infty$. Pruebe que $M=\ker\varphi$ es cerrado y
\[
 \norm{[x]}_{V/M}=\frac{|\varphi(x)|}{C}.
\]
Después considere $V=\ell^1$ y $\varphi(x)=\sum_{n\ge1}\frac{n}{n+1}x_n$. Demuestre que $C=1$ y que la clase cuyos representantes satisfacen $\varphi(x)=1$ tiene norma $1$, pero no contiene ningún vector de norma $1$. Ejercicio propio de esta edición.
\end{ejercicio}


Los cocientes permiten medir distancia a un subespacio y las seminormas obligan a precisar qué vectores se identifican. Volvemos ahora a la dimensión finita: allí las coordenadas proporcionan controles que pueden perderse en los espacios de los últimos ejemplos.

\section{Subespacios de dimensión finita y continuidad}\label{lectura:22}\label{sec:nuevo08:finita}

\begin{proposicion}\label{ch22:suma-cerrada}
Si $M$ es un subespacio cerrado de un espacio normado $V$ y $F$ es un subespacio de dimensión finita, entonces $M+F$ es cerrado.
\end{proposicion}
\begin{proof}
Empecemos por $F=\F u$. Si $u\in M$, la suma es $M$ y no hay nada que probar. Si $u\notin M$, sea $d=\dist(u,M)>0$. Para cualquier escalar $t$, la invariancia de $M$ bajo multiplicación da
\[
 \dist(tu,M)=|t|d;
\]
para $t=0$ ambos lados son cero, y para $t\ne0$ se obtiene cambiando $m$ por $m/t$ en el ínfimo.

Sea $z_n=m_n+t_nu$ una sucesión de $M+\F u$ que converge a $z$ en $V$. Como $m_n-m_k\in M$,
\[
 |t_n-t_k|d\leq\norm{z_n-z_k}.
\]
La convergencia de $(z_n)$ implica que es de Cauchy, de modo que $(t_n)$ es de Cauchy en $\F$ y converge a algún $t$. Entonces $m_n=z_n-t_nu\to z-tu$. Por cerradura de $M$, $z-tu\in M$, y así $z\in M+\F u$. Esto prueba que la suma es cerrada.

Para $F=\spanop\{u_1,\ldots,u_k\}$, aplique el caso anterior sucesivamente a
\[
 M,\quad M+\F u_1,\quad M+\F u_1+\F u_2,\quad\ldots.
\]
Cada paso añade una dirección a un subespacio ya cerrado. Después de un número finito de pasos se obtiene $M+F$ cerrado.
\end{proof}

\begin{corolario}\label{ch22:finito-cerrado}
Todo subespacio de dimensión finita de un espacio normado es cerrado.
\end{corolario}
\begin{proof}
Tome $M=\{0\}$ en la proposición anterior. Esta conclusión coincide con la obtenida por coordenadas en el lema~\ref{ch19:coordenadas}; la nueva prueba resalta la estabilidad de la cerradura al añadir un número finito de direcciones.
\end{proof}



\begin{corolario}\label{ch22:lineal-finito}
Si $V$ tiene dimensión finita, toda transformación lineal $T:V\to W$ es continua. Todas las normas en $V$ son equivalentes y $V$ es de Banach.
\end{corolario}
\index{dimension finita@dimensión finita!equivalencia de normas}
\begin{proof}
Si $V=\{0\}$, todas las conclusiones son inmediatas. Supongamos $\dim V=k>0$ y elija una base $v_1,\ldots,v_k$. Por el lema~\ref{ch19:coordenadas} existe $c>0$ tal que $\norm{\sum_j a_jv_j}\geq c\max_j|a_j|$. Para una transformación lineal $T:V\to W$ y $v=\sum_j a_jv_j$, se obtiene la cota explícita
\[
 \norm{Tv}\leq\sum_{j=1}^k|a_j|\norm{Tv_j}
 \leq\frac{\sum_{j=1}^k\norm{Tv_j}}{c}\norm v=:C_T\norm v.
\]
Por linealidad, $\norm{Tx-Ty}\leq C_T\norm{x-y}$. Si $C_T>0$, una distancia de entrada menor que $\eps/C_T$ produce una distancia de salida menor que $\eps$; si $C_T=0$, $T$ es el operador cero. Esto demuestra continuidad sin utilizar todavía la teoría general de la norma de operador.

Si $\norm{\cdot}_1$ y $\norm{\cdot}_2$ son dos normas en $V$, la identidad de $(V,\norm{\cdot}_1)$ en $(V,\norm{\cdot}_2)$ y la identidad inversa son lineales. Por lo recién probado, ambas son continuas. Aplicando la cota explícita anterior a ambas identidades, existen $c,C>0$ tales que $c\norm v_1\leq\norm v_2\leq C\norm v_1$. Estas desigualdades son la equivalencia de normas.

Elija una base de $V$ y la norma de coordenadas $\norm{\sum a_jv_j}_0=\max_j|a_j|$. Una sucesión de Cauchy para esa norma es de Cauchy en cada coordenada; la completitud de $\F$ produce sus límites y, por tanto, un límite en $V$. Así $\norm{\cdot}_0$ es completa. Las desigualdades de equivalencia muestran que una sucesión de Cauchy en la norma original lo es en $\norm{\cdot}_0$, y que la convergencia en esta implica convergencia en la original. Por tanto $V$ es de Banach.
\end{proof}

\begin{proposicion}[Estabilidad de la independencia]\label{ch22:independencia}
Si $v_1,\ldots,v_k$ son linealmente independientes en $V$, existe $\delta>0$ tal que $\norm{w_j-v_j}<\delta$ para todo $j$ implica que $w_1,\ldots,w_k$ también son linealmente independientes.
\end{proposicion}
\begin{proof}
El lema~\ref{ch19:coordenadas} proporciona $c>0$ con
\[
 \norm{\sum_{j=1}^k a_jv_j}\geq c\max_j|a_j|
 \geq\frac ck\sum_{j=1}^k|a_j|.
\]
Tome $\delta=c/(2k)$. Si $\sum_j a_jw_j=0$, entonces
\begin{align*}
 \frac ck\sum_j|a_j|
 &\leq\norm{\sum_j a_jv_j}
 =\norm{\sum_j a_j(v_j-w_j)}\\
 &\leq\sum_j|a_j|\norm{v_j-w_j}
 \leq\frac c{2k}\sum_j|a_j|.
\end{align*}
Esto fuerza $\sum_j|a_j|=0$, y por ello todos los coeficientes son cero.
\end{proof}

\section{Compacidad y dimensión}

\begin{teorema}[Heine--Borel en dimensión finita]\label{ch22:heine-borel}
En un espacio normado de dimensión finita, un conjunto es compacto si y solo si es cerrado y acotado.
\end{teorema}
\begin{proof}
En todo espacio métrico, los compactos son cerrados y acotados. Para el recíproco, use una base para identificar $V$ linealmente con $\F^k$. El corolario~\ref{ch22:lineal-finito} muestra que la identificación y su inversa son continuas y que las normas se acotan mutuamente. Por tanto la imagen de un cerrado acotado de $V$ es cerrada y acotada en $\F^k$. Heine--Borel euclidiano la hace compacta, y su imagen por la identificación inversa es compacta en $V$.
\end{proof}

\begin{teorema}\label{ch22:caracterizacion}
Para un espacio normado $V$, las siguientes condiciones son equivalentes:
\begin{enumerate}[label=(\roman*)]
\item $V$ tiene dimensión finita;
\item la bola cerrada $\overline B(0,1)$ es compacta;
\item $V$ es localmente compacto.
\end{enumerate}
\end{teorema}
\index{compacidad!local en espacios normados}
\begin{proof}
Si la dimensión es finita, Heine--Borel da la compacidad de todas las bolas cerradas. Cada punto tiene una de esas bolas como vecindad compacta, por lo que (i) implica (ii) y (iii).

Si $V$ es localmente compacto, el origen tiene una vecindad compacta $K$. Existe $r>0$ con $B(0,r)\subseteq K$. La bola cerrada $\overline B(0,r/2)$ es un cerrado de $V$ contenido en $K$, luego es compacta. La homotecia $v\mapsto (2/r)v$ muestra que $\overline B(0,1)$ es compacta. Esto prueba (iii)$\Rightarrow$(ii).

Supongamos (ii). La compacidad implica acotación total, así que la bola unitaria está cubierta por un número finito de bolas abiertas de radio $1/2$, con centros $v_1,\ldots,v_k\in V$. Sea $M=\spanop\{v_1,\ldots,v_k\}$. Es cerrado por ser de dimensión finita. Si $M\ne V$, el lema de Riesz daría $v$ con $\norm v=1$ y $\dist(v,M)>1/2$. Sin embargo, algún centro satisface $\norm{v-v_j}<1/2$, contradicción. Por tanto $M=V$ y $V$ tiene dimensión finita.
\end{proof}

\begin{corolario}\label{ch22:no-sigma-compacto}
Un espacio de Banach infinito dimensional no es unión numerable de conjuntos compactos.
\end{corolario}
\begin{proof}
Si $V=\bigcup_nK_n$ con cada $K_n$ compacto, todos son cerrados. Baire da un $K_n$ con interior no vacío. Existe una bola $B(x,r)\subseteq K_n$. La bola cerrada $\overline B(x,r/2)$ es un cerrado contenido en el compacto $K_n$, y por ello es compacta. Trasladar y reescalar hace compacta la bola unitaria de $V$. El teorema anterior contradice la dimensión infinita.
\end{proof}

\begin{ejemplo}
En $\ell^2$, los vectores canónicos $e_n$ tienen norma $1$ y $\norm{e_n-e_m}=\sqrt2$ si $n\ne m$. La sucesión no tiene ninguna subsucesión de Cauchy y, por tanto, ninguna convergente. La bola unitaria es cerrada y acotada, pero no compacta. El teorema anterior demuestra que siempre existe una obstrucción de este tipo en dimensión infinita, aunque no haya una base canónica disponible.
\end{ejemplo}

\section{Ejercicios propuestos: operaciones, series y subespacios}

\begin{ejercicio}[Lectura 19, ejercicio 1]\label{ej:lectura-19-1}
Para $u,v\ne0$, ponga $u'=u/\norm u$ y $v'=v/\norm v$. Demuestre que
\[\norm{u'-v'}\leq\frac{2\norm{u-v}}{\norm u}.\]
\end{ejercicio}


\begin{ejercicio}[Lectura 19, ejercicio 2]\label{ej:lectura-19-2}
Para $A,B\subseteq V$, sea $A+B=\{a+b:a\in A,b\in B\}$. Pruebe: (i) si $A$ es abierto, $A+B$ es abierto; (ii) si $A,B$ son compactos, $A+B$ es compacto; (iii) si $A$ es cerrado y $B$ compacto, $A+B$ es cerrado. Dé además dos conjuntos cerrados cuya suma no sea cerrada.
\end{ejercicio}


\begin{ejercicio}[Lectura 19, ejercicio 3]\label{ej:lectura-19-3}
Demuestre que todo subespacio propio $M$ de $V$ tiene interior vacío.
\end{ejercicio}


\begin{ejercicio}[Lectura 19, ejercicio 4]\label{ej:lectura-19-4}
Pruebe que $V$ es separable si y solo si $\overline B(0,1)$ es separable.
\end{ejercicio}


\begin{ejercicio}[Lectura 19, ejercicio 5]\label{ej:lectura-19-5}
Si $(v_n)$ es de Cauchy, $v_n\ne0$, y no converge a $0$, pruebe que $(v_n/\norm{v_n})$ es de Cauchy. Deduzca que $V$ es de Banach si y solo si $\overline B(0,1)$ es completa.
\end{ejercicio}


\begin{ejercicio}[Lectura 19, ejercicio 6]\label{ej:lectura-19-6}
Sean $M_1,M_2\leq V$ tales que $\cl{M_1}\cap M_2=\{0\}$ y $M_1+M_2$ es cerrado. Demuestre que $M_1$ es cerrado.
\end{ejercicio}


\begin{ejercicio}[Lectura 19, ejercicio 7]\label{ej:lectura-19-7}
Si $v_n\to v$, demuestre que las medias de Cesàro $u_n=(v_1+\cdots+v_n)/n$ convergen a $v$.
\end{ejercicio}


\section{Ejercicios propuestos: dimensión y mejores aproximantes}

\begin{ejercicio}[Lectura 22, ejercicio 4]\label{ej:lectura-22-4}
Si $M$ tiene dimensión finita y $v\notin M$, demuestre que existe $m_0\in M$ con $\dist(v,M)=\norm{v-m_0}$. Deduzca que existe $u$ con $\norm u=\dist(u,M)=1$.
\end{ejercicio}


\begin{ejercicio}[Lectura 22, ejercicio 5]\label{ej:lectura-22-5}
Sea $c_{00}=\spanop\{e_n:n\in\N\}$, el espacio de sucesiones de soporte finito. Demuestre que no es completo con ninguna norma.
\end{ejercicio}


En dimensión finita, las coordenadas controlan continuidad, completitud y compacidad. El capítulo siguiente estudia cómo cambian estas propiedades al considerar operadores entre espacios normados, especialmente cuando la dimensión es infinita.
